\documentclass[10pt]{article}

% =====================================================================
%  SUPPLEMENTARY INFORMATION
%  for: "Data extraction from published articles on atomic spectroscopy
%        with large language models -- A proof of concept based on the
%        Pr^{2+} spectrum"
%
%  This is a STANDALONE document. Nature journals post the Supplementary
%  Information as a separate PDF that is NOT copy-edited or typeset by
%  the publisher, and it carries its OWN, independent reference list.
%  Compile it on its own in Overleaf (biblatex + biber), exactly like
%  main.tex. Do NOT \input this file into main.tex.
% =====================================================================

% --- PACKAGES (mirror of main.tex) ---
\usepackage[utf8]{inputenc}
\usepackage[T1]{fontenc}
\usepackage{amsmath}
\usepackage{amssymb}
\usepackage{graphicx}
\usepackage{xurl}
\usepackage{hyperref}
\usepackage{enumitem}
\usepackage{geometry}
\usepackage{authblk}
\usepackage{pgfplots}
\usepackage{multirow}
\usepackage{booktabs}      % nicer rules for the data-dictionary tables
\usepackage{longtable}     % multi-page tables (data dictionaries)
\usepackage{tikz}
\usepackage{array}
\usepackage{nicefrac}
\usepackage[backend=biber, style=numeric]{biblatex}
\setcounter{biburllcpenalty}{100}
\setcounter{biburlucpenalty}{100}
\setcounter{biburlnumpenalty}{100}
\usepackage{orcidlink}
\usepgfplotslibrary{groupplots}
\pgfplotsset{compat=1.18}

% Author-annotation colours (same as main.tex, so notes survive a copy/paste)
\newcommand{\DB}[1]{{\color{blue}{{#1}}}}
\newcommand{\MP}[1]{{\color{red}{{#1}}}}
\newcommand{\YR}[1]{{\color{teal}{{#1}}}}
\newcommand{\AK}[1]{{\color{brown}{{#1}}}}

% --- SUPPLEMENTARY NUMBERING (Nature convention) ---
% Figures -> "Supplementary Fig. 1", Tables -> "Supplementary Table 1",
% equations -> "(S1)". Sections are numbered as "Supplementary Note N".
\renewcommand{\thefigure}{S\arabic{figure}}
\renewcommand{\thetable}{S\arabic{table}}
\renewcommand{\theequation}{S\arabic{equation}}
\renewcommand{\figurename}{Supplementary Fig.}
\renewcommand{\tablename}{Supplementary Table}

% Two independent .bib sources -> a single SI reference list numbered from 1.
\addbibresource{references.bib}
\addbibresource{references_supp.bib}

\title{Supplementary Information\\[4pt]
\large Data extraction from published articles on atomic spectroscopy with
large language models --- A proof of concept based on the Pr$^{2+}$ spectrum}

\author[1]{Manas Pandey}
\author[2,3,4]{Alexander Kramida}
\author[5,6,7,8,*]{Dmitry Budker}
\author[2,3,4]{Yuri Ralchenko}
\author[2,3,4]{Rodrigo Ibacache}
\affil[ ]{\footnotesize Affiliations as in the main article. $^{*}$budker@uni-mainz.de}

\begin{document}
\maketitle

% =====================================================================
\noindent
This Supplementary Information documents the second part of the project
(``LineClass''): the classification of the observed spectral lines of
Pr~III and the refinement of its energy levels. It provides (i) the
preparation of the input data (Supplementary Note~1), (ii) the
classification and level-refinement algorithm (Supplementary Note~2),
(iii) the statistical validation of the results (Supplementary Note~3),
and (iv) the description (data dictionaries) of the two large machine-readable
data tables released with the article (Supplementary Note~4). The reference
implementation is the Python code \texttt{classify\_lines.py}.
% =====================================================================


% #####################################################################
\section*{Supplementary Note 1: Preparation of the input data}
\addcontentsline{toc}{section}{Supplementary Note 1: Preparation of the input data}
\label{sn:dataprep}
% #####################################################################

The classification algorithm takes three inputs: a complete list of known
energy levels, a list of observed lines with estimated wavenumber
uncertainties and intensities reduced to a common scale, and a list of
calculated transition intensities with assigned uncertainties. Each was
prepared as follows.

\subsection*{S1.1 Energy-level compilation}
\label{ssec:levels}
All previously known Pr~III levels were assembled by
merging (a) the levels adopted in the NIST Atomic Spectra
Database~\cite{NIST_ASD} (Sugar's levels~\cite{Sugar1969} as
revised by Martin \textit{et~al.}~\cite{Martin1978}) with (b) the
unpublished set of new and revised levels constructed by J.-F.~Wyart,
communicated to us through P.~Palmeri~\cite{Wyart_unpub}. Each level
carries a unique identifier, energy (cm$^{-1}$), parity, and total
angular momentum $J$. Wyart had revised two energy level values compared to NIST ASD.
His list had limited labeling information and no eigenvector composition. 
Thus, for all levels in common, we used data from NIST ASD. The total number 
of levels in the combined list is 593.

\subsection*{S1.2 Assignment of wavenumber uncertainties}
\label{ssec:unc}
Each observed wavenumber was assigned an estimated $1\sigma$ uncertainty
based on Sugar's general statements on measurement precision~\cite{Sugar1969,Sugar1974}
together with line character and agreement of the observed wavelength with the Ritz value
calculated from Sugar's energy levels. These
uncertainties set the matching tolerances used throughout the algorithm.

Sugar~\cite{Sugar1974} states that his wavelengths measured in air
``have an estimated accuracy of \~0.003~{\AA} except for lines that are 
blended or complex, usually as a result of hyperfine structure (hfs), which 
show an average deviation of $\pm$0.007~{\AA} from predicted values.''
Thus, for most of the air wavelengths given with no character, we adopted the uncertainty 
of 0.003~{\AA}. This agrees fairly well with the root-mean-square difference (rms) 
from Sugar's energy levels for these lines, rms$_{\mathrm{O}-\mathrm{R}} = 0.004$~{\AA}. However, several tens 
of such wavelengths deviate from Sugar's Ritz values by 0.012~{\AA} to 0.054~{\AA}.
Sugar's classifications of these lines make sense only if their measurement 
uncertainties were greater than average. Thus, for these wavelengths we adopted 
uncertainties equal to the absolute value of their difference from Sugar's Ritz values ($\delta\lambda_{\mathrm{O}-\mathrm{R}}$).
For complex lines (character `c' in \cite{Sugar1974}), the rms wavelength difference from 
Sugar's Ritz values is 0.029~{\AA}, which we adopted as their uncertainties.
For all other line characters not involving partially resolved hfs 
(i.e., `h', `ch', `bl', `d', `w', denoting hazy, complex and hazy, blended, double,
and wide lines, respectively), as well as for multiply classified lines, we adopted 
Sugar's nominal uncertainty of 0.007~{\AA}, as it agreed well with rms$_{\mathrm{O}-\mathrm{R}}$.
There were 17 exceptions where $\vert\delta\lambda_{\mathrm{O}-\mathrm{R}}\vert$ was 
exceptionally large. For example, for the line at 8711.375~{\AA} (character `ch'), 
$\delta\lambda_{\mathrm{O}-\mathrm{R}} = -0.40$~{\AA}.
In all these cases, we adopted the uncertainty equal to $\vert\delta\lambda_{\mathrm{O}-\mathrm{R}}\vert$.

The lines marked as the strongest components of a partially resolved hfs were treated differently.
These correspond to line characters `*r' and `*v' (flag patterns shaded to the red and to the violet,
respectively). Reader and Sugar~\cite{Reader1965} measured the total width of hfs patterns for a few 
tens of Pr~III lines in the wavelength range between 2841~{\AA} and 3578~{\AA}. According to their
observations, at least for the subset of energy levels involved, the width of the hfs patterns is
nearly constant in terms of wavenumbers ($\approx$1.5~cm$^{-1}$). Thus, the deviation of the position of 
the strongest component from the center of gravity is expected to be roughly independent 
of wavelength if expressed in wavenumber units. The rms difference of observed wavenumbers from Sugar's
Ritz values was found to be 0.08~cm$^{-1}$ for lines with character `*v' and 0.09~cm$^{-1}$ for character 
`*r'. These values were adopted as wavenumber uncertainties for most of Sugar's classified lines having
these characters. Again, a few lines with exceptionally large deviations from Ritz values have been 
assigned larger uncertainties (up to 1~cm$^{-1}$) to make their assignments consistent with the 
measured wavenumbers. For unclassified lines with these characters, we conservatively adopted an 
uncertainty of 0.5~cm$^{-1}$. For this category of lines, the assigned uncertainties pertain to the
precision with which the measured peak position corresponds to the center of gravity of the
hfs pattern of the fine-structure transition assigned to it, rather than to the measurement precision
of the peak position. For all lines, the described method of uncertainty estimation combines 
measurement precision with interpretation accuracy. For example, the position of a blended line could be 
measured quite precisely, but this measured peak position would not exactly match the Ritz wavelength 
because of the blending.

Uncertainties of Sugar's measurements of 1969~\cite{Sugar1969}, which were made in vacuum, were evaluated
in a similar way. Sugar described them as follows: "From the variation of wavelengths obtained with different 
exposures the accuracy of the final values is estimated to be $\pm 0.004$~{\AA}. Systematic shifts between
the various sets of measurements were removed by standardizing on wavelengths of internal impurities of
C, N, 0, and Si." Based on that, we assumed that there are no systematic errors in Sugar's measurements. 
All lines given in \cite{Sugar1969} without any character were assigned the uncertainty of 0.004~{\AA},
except for four lines that showed a significantly larger $\delta\lambda_{\mathrm{O}-\mathrm{R}}$ (1685.033~{\AA}, 
1931.739~{\AA}, 1933.865~{\AA}, and 2057.554~{\AA}). For these lines, we increased the uncertainties by a 
factor ranging from 4 to 12. For the lines having any non-empty character, we adopted a twice greater 
uncertainty, 0.008~{\AA}. With these adopted uncertainties ($u$), the rms of the normalized differences 
($\delta\lambda_{\mathrm{O}-\mathrm{R}}/u$ is 1.03, which indicates that all Sugar's classifications in
the vacuum region are statistically consistent with the adopted measurement uncertainties.

\subsection*{S1.3 Transition-probability data and their uncertainties}
\label{ssec:tp}
The primary source of calculated transition rates was the Mons DREAM
database~\cite{DREAM}. Because the DREAM list for Pr~III was derived from
most---but not all---of Wyart's energy levels and covers only a limited
wavelength range, its coverage of the full transition space is partial.

An experimental alternative exists. Ryder~\cite{Ryder2011}, as part of a
broader study of several singly-ionized, doubly-ionized, and neutral
lanthanide and transition-element spectra, reported laser-induced breakdown
spectroscopy measurements of a few hundred $A$-values of Pr~III. However, these
measurements suffered from line blending and possibly from
self-absorption and calibration errors. A comparison of the line intensities
observed by Sugar~\cite{Sugar1969,Sugar1974} with the $A$-values of
DREAM~\cite{DREAM} and with those of Ryder~\cite{Ryder2011} reveals that the
average accuracy of Ryder's data is no better than that of DREAM: the
disagreements between Sugar and Ryder amount to a factor of 10 on average and
reach up to a factor of 3000, whereas the DREAM values agree with Sugar's
observed intensities much better, within a factor of 3 on average. 
Furthermore, about half of the Pr~III lines reported by Ryder are absent in 
Sugar's line list, which raises serious doubts about correctness of Ryder's identifications.
For these reasons, Ryder's data were not used in the present work.

The DREAM data were complemented with the theoretical oscillator strengths of
Palmeri~\textit{et al.}~\cite{Palmeri2000} (hereafter P00), who published
calculated $\log(gf)$ values for 215 transitions, of which 185 are in common
with DREAM.
Both DREAM and Palmeri~\textit{et al.} had used a similar method
(Cowan's code with core-polarization corrections). However, the DREAM data 
were updated with the later calculations of Bi\'emont~\textit{et
al.}~\cite{Biemont2001}, which revised the treatment of the core-polarization
effects relative to Palmeri~\textit{et al.}~\cite{Palmeri2000}; therefore, the DREAM
values should be more accurate. We adopted the DREAM
values where available and added only the 30 values from
Palmeri~\textit{et al.}~\cite{Palmeri2000} that are absent from DREAM.

Uncertainties were assigned to every transition rate following the general
method of Kramida~\cite{Kramida2024_EPJD}. The central quantity is the
uncertainty of the calculated line strength $S$ expressed as a function of $S$
itself, $u=u(S)$. To derive it, the 185 transitions common to the two
theoretical datasets were divided by line strength into five sections:
$S>4.9$, $2.28<S\le4.9$, $0.81<S\le2.28$, $0.1<S\le0.81$, and $S\le0.1$~a.u.,
the $S$ values being those of P00. The section boundaries were established by
visual inspection of the left panel of Supplementary Fig.~\ref{fig:tp_unc},
which plots $\log_{10}(f_{\mathrm{P00}}/F_{\mathrm{DREAM}})$ against
$S_{\mathrm{P00}}$ and reveals a consistent pattern of growing discrepancies
with decreasing line strength. For each section, we computed the mean line
strength $\langle S\rangle$ and the rms of
$\log_{10}(f_{\mathrm{P00}}/f_{\mathrm{DREAM}})$. These rms values fall close
to a straight line when plotted against $(\ln\langle S\rangle)^{1/3}$, as shown
in the right panel of the Supplementary Fig.~\ref{fig:tp_unc}. Using the fitted
line, the estimated percentage uncertainty in $S$ was taken as
\begin{equation}
u(S) = \left[\,10^{\,\max\!\left(0.05,\
\left|\log_{10}(f_{\mathrm{P00}}/f_{\mathrm{DREAM}})\right|,\
-0.312925\,(\ln S)^{1/3} + 0.454808\right)} - 1\,\right]\times 100 ,
\label{eq:uS}
\end{equation}
where the power of 10 is the maximum of three quantities: a floor of $0.05$, the lowest adopted
uncertainty of $\log_{10}S$; the actual discrepancy
$\left|\log_{10}(f_{\mathrm{P00}}/f_{\mathrm{DREAM}})\right|$ between the P00
and DREAM values where available; and the value returned by the linear fit of
the right panel of the Supplementary Fig.~\ref{fig:tp_unc}. The same function
$u(S)$ was applied to both datasets, with $S$ taken from the corresponding set
(DREAM or P00). Internally, the line-classification algorithm stores each uncertainty as a
logarithmic quantity $\ln(u_{\%gA}/100 + 1)$.

\begin{figure}[t]
  \centering
  % =====================================================================
%  Two-panel figure for Supplementary Note S1.3
%  (uncertainties of the calculated transition rates).
%
%  Requires: pgfplots + \usepgfplotslibrary{groupplots} and the data file
%  fig_tp_leftpanel.dat in the same directory (both already set up in
%  supplementary.tex). \input this file inside a figure environment.
%
%  Left panel : log10(f_P00 / f_DREAM) vs S_P00 for the 185 transitions
%               common to P00 and DREAM, with the five section boundaries
%               (S = 0.1, 0.81, 2.28, 4.9 a.u.) shown as dotted lines.
%  Right panel: per-section rms of log10(f_P00 / f_DREAM) vs (ln<S>)^{1/3},
%               the fitted straight line, and the 0.05 floor used in u(S).
% =====================================================================
\begin{tikzpicture}
\begin{groupplot}[
    group style={group size=2 by 1, horizontal sep=1.7cm},
    width=6.2cm, height=5.9cm,
    tick align=outside,
    tick label style={font=\footnotesize},
    label style={font=\small},
    legend style={font=\footnotesize, draw=none, fill=none},
]
% ------------------------------ LEFT PANEL ------------------------------
\nextgroupplot[
    xmode=log,
    xlabel={$S_{\mathrm{P00}}$ (a.u.)},
    ylabel={$\log_{10}(f_{\mathrm{P00}}/f_{\mathrm{DREAM}})$},
    xmin=0.006, xmax=70, ymin=-1.3, ymax=1.2,
    ymajorgrids, grid style={gray!20},
]
\addplot[only marks, mark=*, mark size=0.9pt, color=blue!55!black, opacity=0.75]
    table[x=S, y=dlgf] {fig_tp_leftpanel.dat};
\addplot[black, dashed, line width=0.5pt] coordinates {(0.006,0) (70,0)};
% section boundaries S = 0.1, 0.81, 2.28, 4.9 a.u.
\draw[gray, densely dotted] (axis cs:0.1, -1.3)  -- (axis cs:0.1, 1.2);
\draw[gray, densely dotted] (axis cs:0.81,-1.3)  -- (axis cs:0.81,1.2);
\draw[gray, densely dotted] (axis cs:2.28,-1.3)  -- (axis cs:2.28,1.2);
\draw[gray, densely dotted] (axis cs:4.9, -1.3)  -- (axis cs:4.9, 1.2);
% ------------------------------ RIGHT PANEL -----------------------------
\nextgroupplot[
    xlabel={$(\ln\langle S\rangle)^{1/3}$},
    ylabel={rms$\,[\log_{10}(f_{\mathrm{P00}}/F_{\mathrm{DREAM}})]$},
    xmin=-1.7, xmax=1.6, ymin=0, ymax=1.0,
    ymajorgrids, grid style={gray!20},
    legend pos=north east,
]
% linear fit: rms = -0.312925 (ln<S>)^{1/3} + 0.454808
\addplot[red, line width=1pt, domain=-1.65:1.55] {-0.312925054856973*x + 0.454808404431693};
\addlegendentry{linear fit}
% floor value 0.05 (in log10 S) used in u(S)
\addplot[black, dashed, line width=0.5pt, domain=-1.65:1.55] {0.05};
\addlegendentry{floor $=0.05$}
% five per-section rms points
\addplot[only marks, mark=square*, mark size=2pt, color=blue!55!black] coordinates {
    ( 1.3749445390308193, 0.050077477002810282)
    ( 1.0836407228996769, 0.069208070997677076)
    ( 0.66052253734860689, 0.26659558829203356)
    (-1.0039511592184323, 0.78413270900177245)
    (-1.4458935997216575, 0.89459900325234076)
};
\addlegendentry{section rms}
\end{groupplot}
\end{tikzpicture}

  \caption{Estimation of the uncertainties of the calculated Pr~III transition
  rates from the 185 transitions common to the DREAM database~\cite{DREAM} and
  Palmeri~\textit{et al.} (P00)~\cite{Palmeri2000}.
  \textbf{Left:} the ratio $\log_{10}(f_{\mathrm{P00}}/f_{\mathrm{DREAM}})$ as a
  function of the P00 line strength $S_{\mathrm{P00}}$ (logarithmic axis); the
  discrepancies grow systematically toward small line strengths. The dotted
  vertical lines mark the boundaries of the five line-strength sections
  ($S=0.1,\,0.81,\,2.28,\,4.9$~a.u.) chosen by visual inspection.
  \textbf{Right:} the per-section rms of
  $\log_{10}(f_{\mathrm{P00}}/f_{\mathrm{DREAM}})$ versus
  $(\ln\langle S\rangle)^{1/3}$ (squares), the fitted straight line (red),
  $-0.312925\,(\ln\langle S\rangle)^{1/3}+0.454808$, and the floor of $0.05$
  (dashed) adopted as the lowest uncertainty of $\log_{10}S$. The fitted line
  defines the line-strength uncertainty $u(S)$ of Eq.~\eqref{eq:uS}, applied to
  both datasets.}
  \label{fig:tp_unc}
\end{figure}

\subsection*{S1.4 Reduction of intensities to a common scale}
\label{ssec:intens}
The observed intensities reported by Sugar~\cite{Sugar1969,Sugar1974} were
placed on a common, approximately linear scale using the method first
suggested by Kramida~\cite{Kramida_2013FuST_63_313} and detailed by
Kramida~\textit{et al.}~\cite{Kramida2022_ApJS}. The procedure simultaneously
produces a set of \emph{calculated} intensities on the same scale, obtained
from a simplified Boltzmann model of the level populations combined with the
adopted transition rates of Supplementary Note~S1.3. It rests on two mutually
consistent constructions, shown in the two panels of
Supplementary Fig.~\ref{fig:intens} and iterated to convergence.

\emph{Boltzmann plot} (left panel). For every transition with an adopted $gA$
(Supplementary Note~S1.3; mostly DREAM values, a few from P00), we plot
$\ln(I_{\mathrm{obs.cor}}\,\lambda/gA)$ against the upper-level energy
$E_{\mathrm{up}}$, where $I_{\mathrm{obs.cor}}$ is the corrected observed
intensity defined below and $\lambda$ is the transition (vacuum) wavelength.
On a linear intensity scale these points scatter about a straight line of
slope $-1/kT$; the red dotted linear fit yields an effective excitation
temperature $T\simeq1.6$~eV ($kT\simeq1.27\times10^{4}$~cm$^{-1}$) and an
intercept $C$.

\emph{Inverse response function} (right panel). Sugar's photographic
intensities are neither linear nor free of a wavelength-dependent
instrumental response. To characterize the latter we form the inverse
response function $R=\ln(I_{\mathrm{calc}}/I_{\mathrm{obs}})$, where
$I_{\mathrm{obs}}$ are Sugar's reported intensities multiplied by $1000$ and
\begin{equation}
I_{\mathrm{calc}} = \frac{gA}{\lambda_{\mathrm{vac}}}\,
\exp\!\left(-\frac{E_{\mathrm{up}}}{kT} + C\right)
\label{eq:Icalc}
\end{equation}
is the intensity predicted by the Boltzmann model with the parameters of the
left panel ($kT$ expressed in cm$^{-1}$, the same units as $E_{\mathrm{up}}$).
Plotted against the vacuum wavelength $\lambda_{\mathrm{vac}}$, $R$ varies
smoothly but with distinct behavior in six wavelength regions
(\texttt{reg1}--\texttt{reg6}) that together span Sugar's combined range;
within each region $R(\lambda)$ is represented by a low-order polynomial fit
(degree 1 to 5; red dotted curves in Supplementary Fig.~\ref{fig:intens}). The
smoothed response defines the corrected observed intensities,
$I_{\mathrm{obs.cor}} = I_{\mathrm{obs}}\,\exp[R_{\mathrm{fit}}(\lambda)]$,
which are fed back into the Boltzmann plot. The two constructions are iterated
until convergence (seven cycles). The resulting $I_{\mathrm{obs.cor}}$ are the
observed intensities on the common linear scale, and the corresponding
$I_{\mathrm{calc}}$ are the calculated intensities used throughout the line
classification (Supplementary Note~2). The effective temperature $T=1.6$~eV
also sets the Doppler width used in the grouping stage of the weeding step
(Supplementary Note~S2.5).

\begin{figure}[t]
  \centering
  % =====================================================================
%  Two-panel figure for Supplementary Note S1.4
%  (reduction of the observed intensities to a common scale).
%
%  Layout: two independent tikzpictures wrapped in a single \mbox so the
%  panels can never stack; the right (response) panel is wider than the left
%  (Boltzmann) panel to relieve VUV crowding.
%
%  Fits ON TOP -- robust method: in each panel the DATA are drawn in one axis
%  and the red fit lines in a SECOND, transparent axis overlaid exactly on the
%  first (same limits, scale only axis, anchored to the data axis' south-west
%  corner).  The whole second axis is emitted after the first, so its lines are
%  unconditionally on top of the first axis' markers -- immune to any
%  intra-axis line/marker ordering.
%
%  Requires: pgfplots and the data files in the same directory:
%    fig_s14_boltzmann.dat                            (left panel)
%    fig_s14_resp_reg1.dat ... fig_s14_resp_reg6.dat  (right panel, per region)
% =====================================================================
\mbox{%
% ============================ LEFT PANEL ============================
\begin{tikzpicture}
% --- data axis ---
\begin{axis}[
    name=lftA, scale only axis,
    width=3.3cm, height=4.8cm,
    tick align=outside,
    tick label style={font=\footnotesize},
    label style={font=\small},
    scaled x ticks=false,
    xlabel={$E_{\mathrm{up}}$ (cm$^{-1}$)},
    ylabel={$\ln(I_{\mathrm{obs.cor}}\,\lambda / gA)$},
    xmin=8000, xmax=143000, ymin=-9, ymax=6,
    xtick={50000,100000},
    xticklabel style={/pgf/number format/1000 sep={\,}},
    ymajorgrids, grid style={gray!20},
]
\addplot[only marks, mark=*, mark size=0.5pt, color=blue!55!black, opacity=0.55]
    table[x=E, y=bolt] {fig_s14_boltzmann.dat};
\end{axis}
% --- overlay fit axis (exactly over the data axis) ---
\begin{axis}[
    name=lftB, at={(lftA.south west)}, anchor=south west, scale only axis,
    width=3.3cm, height=4.8cm,
    xmin=8000, xmax=143000, ymin=-9, ymax=6,
    axis lines=none, xtick=\empty, ytick=\empty,
]
% Boltzmann fit (slope -7.86527e-5, intercept 4.95823; T = 1.58 eV)
\addplot[red,   line width=1.3pt] coordinates {(12846.7,3.9478) (138874.9,-5.9649)};
\end{axis}
\end{tikzpicture}%
\hspace{6mm}%
% ============================ RIGHT PANEL ============================
\begin{tikzpicture}
% --- data axis ---
\begin{axis}[
    name=rgtA, scale only axis,
    width=8.2cm, height=4.8cm,
    tick align=outside,
    tick label style={font=\footnotesize},
    label style={font=\small},
    scaled x ticks=false,
    xlabel={$\lambda_{\mathrm{vac}}$ (\AA)},
    ylabel={$R=\ln(I_{\mathrm{calc}}/I_{\mathrm{obs}})$},
    xmin=800, xmax=11000, ymin=-8.5, ymax=6.2,
    xtick={2000,4000,6000,8000,10000},
    xticklabel style={/pgf/number format/1000 sep={\,}},
    ymajorgrids, grid style={gray!20},
    legend columns=3,
    legend style={at={(0.985,0.985)}, anchor=north east, font=\scriptsize,
                  fill=white, draw=gray!50, /tikz/every even column/.append style={column sep=3pt}},
]
\addplot[only marks, mark=*,        mark size=1.0pt, color=blue!70!black,  opacity=0.7] table[x=wl,y=R]{fig_s14_resp_reg1.dat}; \addlegendentry{reg\,1}
\addplot[only marks, mark=square*,  mark size=1.0pt, color=teal,           opacity=0.7] table[x=wl,y=R]{fig_s14_resp_reg2.dat}; \addlegendentry{reg\,2}
\addplot[only marks, mark=triangle*,mark size=1.3pt, color=orange!90!black,opacity=0.7] table[x=wl,y=R]{fig_s14_resp_reg3.dat}; \addlegendentry{reg\,3}
\addplot[only marks, mark=diamond*, mark size=1.3pt, color=black!65,       opacity=0.8] table[x=wl,y=R]{fig_s14_resp_reg4.dat}; \addlegendentry{reg\,4}
\addplot[only marks, mark=pentagon*,mark size=1.3pt, color=brown!80!black, opacity=0.8] table[x=wl,y=R]{fig_s14_resp_reg5.dat}; \addlegendentry{reg\,5}
\addplot[only marks, mark=x,        mark size=1.5pt, color=black,          opacity=0.8] table[x=wl,y=R]{fig_s14_resp_reg6.dat}; \addlegendentry{reg\,6}
\end{axis}
% --- overlay fit axis (exactly over the data axis) ---
\begin{axis}[
    name=rgtB, at={(rgtA.south west)}, anchor=south west, scale only axis,
    width=8.2cm, height=4.8cm,
    xmin=800, xmax=11000, ymin=-8.5, ymax=6.2,
    axis lines=none, xtick=\empty, ytick=\empty,
]
% per-region fits: thin solid red, on top of the data axis
\addplot[red,line width=1.3pt] coordinates {
(821.94,0.9549) (857.87,0.7547) (893.79,0.5544) (929.72,0.3541) (965.64,0.1539) (1001.57,-0.0464) (1037.49,-0.2466) (1073.42,-0.4469) (1109.35,-0.6471) (1145.27,-0.8474) (1181.20,-1.0476) (1217.12,-1.2479) (1253.05,-1.4482) (1288.97,-1.6484) (1324.90,-1.8487) (1360.82,-2.0489) (1396.75,-2.2492) (1432.68,-2.4494) (1468.60,-2.6497) (1504.53,-2.8499) (1522.49,-2.9501)
};
\addplot[red,line width=1.3pt] coordinates {
(1529.85,0.6797) (1573.97,0.5567) (1618.10,0.4338) (1662.22,0.3108) (1706.34,0.1878) (1750.47,0.0648) (1794.59,-0.0582) (1838.71,-0.1811) (1882.84,-0.3041) (1926.96,-0.4271) (1971.08,-0.5501) (2015.21,-0.6731) (2059.33,-0.7960) (2103.46,-0.9190)
};
\addplot[red,line width=1.3pt] coordinates {
(2107.92,4.5748) (2140.69,3.8063) (2173.47,3.1733) (2206.25,2.6620) (2239.02,2.2587) (2271.80,1.9497) (2304.58,1.7212) (2337.35,1.5593) (2370.13,1.4505) (2402.91,1.3808) (2435.68,1.3366) (2468.46,1.3040) (2501.24,1.2694) (2534.01,1.2189) (2566.79,1.1389) (2599.57,1.0154) (2632.34,0.8349) (2665.12,0.5835) (2697.90,0.2474) (2730.67,-0.1870) (2747.06,-0.4454)
};
\addplot[red,line width=1.3pt] coordinates {
(2753.06,-0.1264) (2840.35,0.2840) (2927.63,0.5417) (3014.91,0.6628) (3102.20,0.6630) (3189.48,0.5581) (3276.77,0.3642) (3364.05,0.0970) (3451.33,-0.2276) (3538.62,-0.5937) (3625.90,-0.9854) (3713.18,-1.3869) (3800.47,-1.7823) (3887.75,-2.1558) (3975.03,-2.4914) (4062.32,-2.7732) (4149.60,-2.9855) (4236.88,-3.1124) (4324.17,-3.1379) (4411.45,-3.0463) (4455.09,-2.9515)
};
\addplot[red,line width=1.3pt] coordinates {
(4449.18,-3.9992) (4532.44,-3.8244) (4615.70,-3.5239) (4698.96,-3.1530) (4782.22,-2.7596) (4865.48,-2.3840) (4948.74,-2.0592) (5032.01,-1.8106) (5115.27,-1.6563) (5198.53,-1.6066) (5281.79,-1.6648) (5365.05,-1.8263) (5448.31,-2.0793) (5531.57,-2.4044) (5614.83,-2.7749) (5698.09,-3.1565) (5781.35,-3.5074) (5864.61,-3.7784) (5947.88,-3.9129) (6031.14,-3.8467) (6072.77,-3.7163)
};
\addplot[red,line width=1.3pt] coordinates {
(6054.68,-3.6887) (6293.91,-3.1496) (6533.13,-3.2123) (6772.35,-3.5636) (7011.57,-3.9809) (7250.79,-4.3210) (7490.02,-4.5078) (7729.24,-4.5207) (7968.46,-4.3828) (8207.68,-4.1488) (8446.91,-3.8939) (8686.13,-3.7009) (8925.35,-3.6493) (9164.57,-3.8031) (9403.80,-4.1990) (9643.02,-4.8346) (9882.24,-5.6566) (10121.46,-6.5489) (10360.69,-7.3209) (10599.91,-7.6957) (10719.52,-7.6210)
};
\end{axis}
\end{tikzpicture}%
}

  \caption{Reduction of Sugar's observed Pr~III intensities to a common linear
  scale. \textbf{Left:} Boltzmann plot of
  $\ln(I_{\mathrm{obs.cor}}\,\lambda/gA)$ versus the upper-level energy
  $E_{\mathrm{up}}$ for the transitions with an adopted $gA$. The red solid
  line is the linear fit; its slope corresponds to an effective temperature
  $T\simeq1.6$~eV and its intercept is $C$ (Eq.~\eqref{eq:Icalc}).
  \textbf{Right:} the inverse response function
  $R=\ln(I_{\mathrm{calc}}/I_{\mathrm{obs}})$ versus the vacuum wavelength
  $\lambda_{\mathrm{vac}}$. The six fitting regions
  (\texttt{reg1}--\texttt{reg6}) are distinguished by marker shape and color
  (legend); the first extends down to Sugar's shortest wavelength,
  $821.94$~\AA{} (its fit below ${\sim}1000$~\AA{} is an extrapolation, no $gA$
  values being available there). The red solid curves are the per-region
  polynomial fits
  (degrees 1, 1, 3, 3, 4, 5) that
  define the smoothed response $R_{\mathrm{fit}}(\lambda)$ and hence the
  correction $I_{\mathrm{obs.cor}} = I_{\mathrm{obs}}\exp[R_{\mathrm{fit}}(\lambda)]$.}
  \label{fig:intens}
\end{figure}


% #####################################################################
\section*{Supplementary Note 2: The line-classification algorithm}
\addcontentsline{toc}{section}{Supplementary Note 2: The line-classification algorithm}
\label{sn:algorithm}
% #####################################################################

\subsection*{S2.1 Overview}
The algorithm is fully deterministic and uses no LLM. Given the prepared
inputs it (1)~enumerates every electric-dipole (E1) transition allowed
between known levels within the observed range, (2)~matches these to the
observed lines, (3)~grades each match on wavenumber and intensity
agreement, (4)~resolves conflicts where one transition matches several
lines, (5)~iteratively ``weeds'' spurious assignments using a per-level
intensity calibration, and (6)~approximately refines the level energies.
Steps 2--6 are wrapped in an outer cycle that repeats until the set of
accepted assignments and the level energies are both stable (largest
first-iteration energy change $< 0.001$~cm$^{-1}$; at most 10 cycles).
Supplementary Fig.~\ref{fig:flow} shows the data flow.

\begin{figure}[t]
  \centering
  % \includegraphics[width=0.9\linewidth]{figS_pipeline.pdf}
  \fbox{\parbox[c][4.5cm][c]{0.9\linewidth}{\centering\MP{[Pipeline / data-flow
  schematic: read $\rightarrow$ generate $\rightarrow$ match \& grade
  $\rightarrow$ resolve $\rightarrow$ weed $\rightarrow$ optimize $\rightarrow$
  output. Adapt the ASCII diagram in the LineClass README.]}}}
  \caption{Data flow of the \texttt{classify\_lines.py} pipeline.}
  \label{fig:flow}
\end{figure}

\subsection*{S2.2 Transition-space generation}
\label{ssec:space}
All level pairs satisfying the E1 selection rules are enumerated:
opposite parity, $|\Delta J| \le 1$, and not $J=0 \leftrightarrow J=0$.
Only pairs whose Ritz wavenumber
$\sigma_{\text{Ritz}} = E_{\text{upper}} - E_{\text{lower}}$ falls in the
observed range $[9327.0,\,121665.0]$~cm$^{-1}$ are retained. Each candidate
inherits its calculated intensity and uncertainty from the transition-rate
list when available, otherwise these are left undefined. The list is
sorted by $\sigma_{\text{Ritz}}$ to enable binary search during matching.

\subsection*{S2.3 Matching and two-dimensional grading}
\label{ssec:grade}
For each observed line, a bisection over the sorted candidate wavenumbers
selects all transitions within $5.5\,\sigma$ of the observed wavenumber.
Every match receives a two-character grade: a \emph{tier} for wavenumber
agreement and a \emph{subgrade} for intensity consistency
(Supplementary Table~\ref{tab:grades}). A match is flagged
\texttt{new}$=0$ if it reproduces an original Sugar classification,
otherwise \texttt{new}$=1$.

\begin{table}[t]
\centering
\caption{Two-dimensional grade. Tier encodes wavenumber agreement
($w=|\sigma_{\text{obs}}-\sigma_{\text{Ritz}}|$); subgrade encodes
intensity consistency, with $u=u_{\text{calc}}$ (log-scale uncertainty)
and $d=|\ln(I_{\text{calc}}/I_{\text{obs}})|-u$. $\ln2$ and $\ln5$ are
the natural logs of 2 and 5.}
\label{tab:grades}
\begin{tabular}{@{}cl@{\hspace{2.5em}}cl@{}}
\toprule
Tier & Wavenumber condition & Subgrade & Intensity condition \\
\midrule
2 & $w \le 2.0\,\sigma$ & A & $d \le 0$ and $u \le \ln2$ \\
3 & $w \le 3.0\,\sigma$ & B & $d \le \ln2$ and $u \le \ln5$ \\
4 & $w \le 4.0\,\sigma$ & C & $\ln2 < d \le \ln5$ and $\ln2 < u \le \ln5$ \\
5 & $w > 4.0\,\sigma$   & D & $d \le \ln5 < u$ \\
  &                     & E & $d > \ln5$ and $u > \ln5$ \\
  &                     & G & no calculated intensity (or none of A--E) \\
\bottomrule
\end{tabular}
\end{table}

\subsection*{S2.4 Conflict resolution}
\label{ssec:conflict}
When one transition matches several observed lines, the candidates are
ranked by the conservative key
$(\texttt{is\_new},\, |\sigma_{\text{obs}}-\sigma_{\text{Ritz}}|,\,
\text{tier},\, z,\, \Delta I)$, where $z$ is the wavenumber residual in
units of $\sigma$ and $\Delta I$ the intensity mismatch. Legacy
assignments are thus favored, then the smallest wavenumber mismatch, and
so on. The retained match is flagged \texttt{F} (conflicting); if an
original assignment was overridden, it is additionally flagged \texttt{R}
(revised). \MP{[Three alternative scoring rules (lexicographic, $\chi^2$,
rank-based) are present in the code but disabled; mention this only if
relevant.]}

\subsection*{S2.5 Iterative intensity weeding}
\label{ssec:weed}
Weeding decides which candidates are \emph{accepted}. It runs an outer
factor-refinement loop around a per-line three-step decision:
\begin{description}[leftmargin=1.4em,style=nextline]
  \item[Per-level intensity factors.] For every level, a factor is
    computed from its accepted, unblended transitions when at least
    $\text{min\_n}=5$ qualifiers exist, as the Mandel--Paule consensus
    weighted mean~\cite{Paule_Mandel_1982} of $\ln(I_{\text{obs}}/I_{\text{calc}})$,
    clamped to $\pm 5$ ln-units. Two factors are kept per level (as the
    upper vs.\ the lower level of a transition), capturing population
    deviations, self-absorption, and systematic errors in the calculated
    rates. Factors are applied with damping ($\alpha=0.5$) as
    $I_{\text{adj}} = I_{\text{orig}}\exp(f_{\text{from}}+f_{\text{to}})$,
    with uncertainties added in quadrature; the original intensities are
    never mutated. All decisions are then reset and the line is re-weeded.
  \item[Step 1 --- clear-cut decisions.] New candidates with Ritz mismatch
    $> 3\sigma$, or contributing $< 10\%$ of the accepted cumulative
    intensity, are rejected; statistically insignificant legacy candidates
    are rejected around $4\%$/$2\%$ thresholds. Survivors are tested by an
    asymmetric intensity $z$-score (looser when the prediction is too
    strong or for resonance lower levels, stricter when too weak).
  \item[Step 2 --- grouping.] Undecided candidates are tested for pair and
    triple acceptance: a group passes if its intensity-weighted centre of
    gravity matches the observed wavenumber and its effective spread---the
    Doppler width ($T=1.6$~eV, $M=140.90765$~u) convolved with the
    measurement $\sigma$---is consistent with the line profile (thresholds
    relaxed for broadened characters \texttt{h}/\texttt{w}).
  \item[Step 3 --- defaults.] Any remaining undecided new candidate is
    rejected; any remaining legacy candidate is accepted. Transitions whose
    acceptance oscillates between passes are blacklisted and fall through
    to this default.
\end{description}
Convergence is reached when no assignment changes acceptance between passes.

\subsection*{S2.6 Approximate energy-level optimization}
\label{ssec:opt}
Each non-ground level is re-estimated as the weighted mean of the energies
implied by its accepted transitions ($E=E_{\text{lower}}+\sigma_{\text{obs}}$
or $E=E_{\text{upper}}-\sigma_{\text{obs}}$). The weight of a component is
$w = \mathrm{BF}^2/\sigma_{\text{wn}}^2$, where the branching fraction
$\mathrm{BF}=I_{\text{calc}}/\sum I_{\text{calc}}$ shares a blend's weight
among its components (LOPT ``centroid'' weighting). This step is
intentionally minimalistic---it only sharpens the Ritz wavenumbers that
drive the assignment decisions. The \emph{final} level optimization is
performed by the program LOPT~\cite{Kramida2011_LOPT}, fed by the accepted
rows and their branching fractions.


% #####################################################################
\section*{Supplementary Note 3: Statistical validation of the results}
\addcontentsline{toc}{section}{Supplementary Note 3: Statistical validation}
\label{sn:validation}
% #####################################################################

\MP{[This is where the statistical evidence that the classifications and
levels are reliable belongs. Suggested content, each with a figure or
table:]}
\begin{itemize}
  \item Distribution of accepted assignments over the grade matrix
        (tier $\times$ subgrade); fraction of new vs.\ legacy.
  \item Histogram of normalized wavenumber residuals
        $(\sigma_{\text{obs}}-\sigma_{\text{Ritz}})/\sigma$ and its
        reduced $\chi^2$; check against the expected unit-variance normal.
  \item Comparison of refined level energies against the input (NIST/Wyart)
        values: distribution of shifts and their significance.
  \item Number of lines now defining each level, before vs.\ after; levels
        that gained/lost supporting lines.
  \item Consistency of the per-level intensity factors, and agreement of
        observed vs.\ calibrated calculated intensities.
  \item Summary counts reconciled with the main text: 593 levels, 7118
        transitions on 4619 observed lines, 5683 transition probabilities,
        1625 of 6244 lines unclassified, $\sim$2000 new classifications.
\end{itemize}

\begin{figure}[t]
  \centering
  \fbox{\parbox[c][4cm][c]{0.9\linewidth}{\centering\MP{[Validation figure(s):
  residual histogram / grade distribution / energy-shift distribution.]}}}
  \caption{Statistical validation of the Pr~III classifications.}
  \label{fig:validation}
\end{figure}


% #####################################################################
\section*{Supplementary Note 4: The energy-level and line-list data tables}
\addcontentsline{toc}{section}{Supplementary Note 4: Data tables}
\label{sn:tables}
% #####################################################################

The two principal results are too large to typeset here and are released as
machine-readable files (see \emph{Data availability} in the main article).
Supplementary Table~\ref{tab:levels_dict} and
Supplementary Table~\ref{tab:lines_dict} define their columns.

\paragraph{Supplementary Data~1 --- Energy levels} ($593$ rows). \MP{[Fill in the
actual column list once the level table is finalized; the entries below are
placeholders based on a standard critically-evaluated level list.]}

\begin{longtable}{@{}p{0.20\linewidth} p{0.72\linewidth}@{}}
\caption{Column dictionary for Supplementary Data~1 (Pr~III energy levels).}
\label{tab:levels_dict}\\
\toprule
Column & Description \\
\midrule
\endfirsthead
\toprule Column & Description \\ \midrule \endhead
\bottomrule \endfoot
\texttt{level\_id}   & Unique level identifier (matches the line list). \\
\texttt{energy}      & Optimized level energy (cm$^{-1}$). \\
\texttt{unc\_energy} & $1\sigma$ uncertainty of the energy (cm$^{-1}$). \\
\texttt{J}           & Total angular momentum. \\
\texttt{parity}      & Parity (\texttt{e}/\texttt{o}). \\
\texttt{config}      & Electronic configuration. \MP{[if available]} \\
\texttt{term}        & Term / leading-percentage designation. \MP{[if available]} \\
\texttt{n\_lines}    & Number of observed lines defining the level. \\
\texttt{source}      & Origin of the level (NIST/Sugar, Wyart, revised). \\
\MP{[\ldots]}        & \MP{[remaining columns to $\sim$15 total]} \\
\end{longtable}

\paragraph{Supplementary Data~2 --- Classified line list} ($\sim$9000 rows,
21 columns, as written by \texttt{classify\_lines.py}).

\begin{longtable}{@{}p{0.20\linewidth} p{0.72\linewidth}@{}}
\caption{Column dictionary for Supplementary Data~2 (Pr~III classified lines).}
\label{tab:lines_dict}\\
\toprule
Column & Description \\
\midrule
\endfirsthead
\toprule Column & Description \\ \midrule \endhead
\bottomrule \endfoot
\texttt{wn\_obs}           & Observed wavenumber (cm$^{-1}$). \\
\texttt{unc\_wn\_obs}      & Wavenumber uncertainty $\sigma$ (cm$^{-1}$). \\
\texttt{obs\_intens}       & Observed intensity, reduced to a uniform scale. \\
\texttt{char}              & Line character (\texttt{h}, \texttt{w}, \texttt{bl}, \ldots). \\
\texttt{low\_id}           & Lower energy-level id. \\
\texttt{upp\_id}           & Upper energy-level id. \\
\texttt{calc\_intens}      & Calculated intensity after per-level factors. \\
\texttt{orig\_calc\_intens}& Original theoretical intensity. \\
\texttt{u\_calc}           & Log-uncertainty of \texttt{calc\_intens}. \\
\texttt{intens\_from\_f}   & Upper-level intensity correction factor (ln). \\
\texttt{intens\_to\_f}     & Lower-level intensity correction factor (ln). \\
\texttt{dif\_wn\_O-C}      & Observed $-$ Ritz wavenumber (cm$^{-1}$). \\
\texttt{grade}             & Two-dimensional grade (tier $2$--$5$ $+$ subgrade A--E/G). \\
\texttt{notes1}            & Conflict flags: \texttt{F} (conflicting), \texttt{R} (revised). \\
\texttt{notes2}            & Per-transition decision trace from weeding. \\
\texttt{new}               & $1$ new, $0$ original Sugar, blank unclassified. \\
\texttt{accepted}          & $1$ accepted, $0$ rejected, blank unclassified. \\
\texttt{low\_E}            & Lower-level energy (cm$^{-1}$). \\
\texttt{upp\_E}            & Upper-level energy (cm$^{-1}$). \\
\texttt{rwn}               & Ritz wavenumber $=\texttt{upp\_E}-\texttt{low\_E}$ (cm$^{-1}$). \\
\texttt{BF}                & Branching fraction of the component (LOPT input; $0$ if not accepted). \\
\end{longtable}

A short excerpt of the line list (the strongest few lines) may be shown
here for illustration:
\MP{[optional excerpt table --- 5--10 representative rows.]}


% #####################################################################
% Independent Supplementary reference list (numbered from 1).
% #####################################################################
\printbibliography[title={Supplementary References}]

\end{document}
